\section{What this article is, and how to read it}\label{sec:reading}

\paragraph{Object.} The object is the Lean~4 source tree \cite{alpotree} (Lean v4.33.1, Mathlib at
commit \texttt{0df444a3}), identified by the SHA-256 hash of its archive. It has 175{,}187
\texttt{.lean} files in ten libraries (Appendix~\ref{app:libraries}); the root file
\leanlocx{src/Main.lean}{1} exports five theorems. The accompanying note \cite{remark} gives a
timeline of twenty ideas (``items 1--20''), which \S\ref{sec:remark} maps to the source.

\paragraph{Status labels.} Every mathematical claim carries one of four labels.
\begin{description}[font=\normalfont]
\item[\sttag{classical}] a published theorem or a standard fact, with a citation.
\item[\sttag{Lean: file:line}] identifies a statement or definition in the source at this
  location. Where we discuss a proof body, the label refers to the body as written. The build,
  axiom check and kernel replay are in progress (\S\ref{sec:verification}).
\item[\sttag{remark, item $n$}] the author's own description in the note.
\item[\sttag{interpretation}] our reading: what a Lean statement means, or which published theorem
  it corresponds to. Lean declaration names are suggestive (the source has names such as
  \texttt{caraianiLocalGlobalDual} and \texttt{Green}) but they are not evidence; our
  interpretations rest on the quoted statements, and remain interpretations.
\end{description}
File paths are relative to the tree's \texttt{src/} directory and line numbers refer to the archive
with the hash above.

\paragraph{Quotations.} Lean statements are printed in framed boxes headed ``Source: file:line''.
Their text is produced by a script from the source, not typed: a theorem is quoted up to its
\texttt{:=} (the proof is elided as \texttt{:= …}), and a definition or structure is quoted with its
body. A checking script (Appendix~\ref{app:method}) regenerates every quotation and fails if any
box differs from the source, if a cited line does not start a declaration, or if a declaration name
in the text does not exist in the tree.

\paragraph{What the article does not do.} It does not judge whether the proof is correct; the
kernel check decides that. It does not trace the source to the bottom: the map follows the proof
from \texttt{Main.lean} downward for six to seven levels, until each branch reaches a named
theorem whose statement can be read on its own. The article compares the input statements reached
on those branches with results from the literature.

\paragraph{Notation.} $\chi$ and $\psi$ are Dirichlet characters; $L(s,\chi)$ is Mathlib's
\mathlib{DirichletCharacter.LFunction}, the meromorphic continuation of
$\sum_{n\ge1}\chi(n)n^{-s}$; $\Lambda$ is the von Mangoldt function; $\gamma$ is Euler's constant.
For a number field $E$, $\mathcal O_E$ is its ring of integers; for a CM field $E$, $E^+$ is its
maximal totally real subfield.
